\section{引言}\label{sec:introduction}

符号邻接矩阵为谱图论、切换理论与离散磁性思想提供了一个具体交汇点。顶点
切换等价于用对角符号矩阵共轭符号邻接矩阵，因而谱只依赖于圈通量，而不依赖
边规范。在一张固定图的所有赋号中最小化谱半径，遂成为有限通量空间上的优化
问题。这一问题既属于经典符号图理论 \cite{Zaslavsky1982}，也与 Bilu--Linial
谱赋号问题 \cite{BiluLinial2006} 相接。后者包括二分 Ramanujan 提升的交错族
方法及显式近 Ramanujan 构造
\cite{MarcusSpielmanSrivastava2015,MohantyODonnellParedes2022}，但这些结果
并不确定固定非二分图上的双侧谱最优值。符号图类中的极值谱半径问题构成另一条
相关研究线索 \cite{BelardoCioabaKoolenWang2019,BelardoBrunettiCiampella2021}。

本文研究四正则循环图

\begin{equation}
G_n=C_n(1,2),
\end{equation}

其中顶点是模 \(n\) 的剩余类，循环距离为一或二的顶点之间连边。这个模型并非
普适模型，但具有三项相互配合的特点：它是含有重叠三角形与四边形的最简单
四正则循环图之一；切换类可用显式圈坐标完整表示；周期赋号可约化为有限维
Floquet 矩阵。更关键的是，算子平方后出现局部因子 \(1+Q_i\)：负四边形通量
消去耦合，正通量激活缺陷。由此，全局谱优化可与局部缺陷几何联系起来，同时
仍保有可完整枚举的有限商空间。

Suvagiya 最近给出了一个三角形通量交替的特殊族
\cite{Suvagiya2026Signed}。当 \(n\) 为
偶数时，该族中 Hamilton 步长一圈和乐为负的两个成员具有谱半径

\begin{equation}
\rho_-(n)=2 \sqrt{\cos^{2}(\frac{\pi}{n})+\cos^{2}(\frac{2\pi}{n})}.
\end{equation}

文献 \cite{Suvagiya2026Signed} 的猜想 3 断言，\(G_n\) 的任何赋号都不可能
具有更小的谱半径。该文验证了 \(n=8,10,\ldots,18\) 的情形，并将一般下界
留作公开问题。

循环矩阵的 Fourier 分析 \cite{Davis1979} 与周期图算子的 Floquet 理论
\cite{Kuchment1993,KorotyaevSaburova2017} 提供算子背景。周期图上的磁通量迹
公式与本文的闭游走矩计算尤为接近 \cite{KorotyaevSaburova2023}。本文所用
“通量相”一词来自磁性算子与物理学先例，包括 Lieb 的定理 \cite{Lieb1994}；
但本文不援引该定理的物理基态原理或假设。

继承自既有工作的要素包括交替三角形族、它的四个切换/和乐类、上述
\(\rho_-(n)\) 公式以及猜想本身 \cite{Suvagiya2026Signed}；配套的奇偶族
框架见 \cite{Suvagiya2026Parity}。本文的贡献并不限于给出一个显式反例，还包括
确定最小失效阶、推广为无限算术族、计算精确谱边、解释局部消去机制、分类全部
周期八相位、导出一般周期必要条件，并完成一个有界周期范围的精确分类。

该猜想是错误的，但其失效出现得很晚：所有小于 32 的允许阶数都满足所提
下界。最小反例是一个 32 阶赋号，其三角形通量字是下列八元组的四次重复：

\begin{equation}
(+,+,-,+,-,-,+,-).
\end{equation}

同一个字定义了一个周期八算子。它的 Floquet 多项式可以精确计算，而一个
一致正性论证将这一显式反例推广为无限族，覆盖所有不小于 32 的八的倍数。
在平方算子中，\(Q_i=-1\) 的项发生消去，而 \(Q_i=+1\) 的项成为活跃缺陷；
该反例在八点胞中把两个正缺陷布置为对径位置。这个机制而非孤立的 32 阶比较，
构成全文的组织主线。

\subsection{主要结果}

第一个结果将一个显式反例与有限穷举排除相结合。

\begin{theorem}[最小反例定理]
\label{thm:smallest-counterexample}
对于每个满足 \(8\leq n\leq 30\) 的偶数 \(n\)，猜想下界成立；而在
\(n=32\) 时失效。因此，32 是最小反例阶数。

\end{theorem}
直至 30 阶的排除是一条有限计算机辅助定理。其证明枚举切换类的一个完整
商集，并用精确有理不等式认证每个非最优代表元。32 阶反例本身则有一个简短的
精确正定性证书。

该显式构造可以周期延拓。

\begin{theorem}[无限反例族定理]
\label{thm:infinite-counterexample-family}
对于每个 \(n=8L\)、\(L\geq 4\)，以及任一 Hamilton 和乐，第~
\ref{sec:periodic-construction} 节定义的周期八赋号都满足

\begin{equation}
\rho(A)^{2} < \frac{1561}{200} < \rho_-(n)^{2}.
\end{equation}

因此，猜想在每个不小于 32 的八的倍数处失效。上述结论不意味着猜想在每个
大于 32 的偶数阶都失效。

\end{theorem}
在无限体积极限下，Floquet 分析给出精确谱边。

\begin{theorem}[周期八精确谱边定理]
\label{thm:exact-period-eight-edge}
对于目标周期八相位，无限体积谱半径的平方为

\begin{equation}
\eta=4+\sqrt{10+2\sqrt{5}},
\end{equation}

并且最高谱带仅在 Bloch 参数 \(z=1\) 处达到 \(\eta\)。

\end{theorem}
其局部机制由一套完整相位分类刻画。对于合法的周期八通量字 \(Q\)，令
\(D(Q)\) 为其中正号位置的集合，并令 \(R(Q)\) 表示无限体积谱半径的平方。

\begin{theorem}[周期八相位三分定理]
\label{thm:eight-barrier-trichotomy}
下列三种情形恰有一种发生：

\begin{equation}
\begin{cases}
R(Q)=8, & D(Q)=\varnothing; \\
R(Q)=\eta<8, & D(Q)=\{j,j+4\}; \\
R(Q)>8, & \text{对于其他每个合法周期八 }Q
\end{cases}
\text{。}
\end{equation}

特别地，在平移、反射和全局边符号取反意义下，目标相位是唯一的周期八最优
相位，而全负相位是唯一的次优相位。

\end{theorem}
同一个平方算子计算还能给出任意周期的信息。令 \(d=\lvert D(Q)\rvert\)，
令 \(a\) 计数相邻正号对，令 \(b\) 计数循环距离为二的正号对。

\begin{theorem}[一般周期必要条件定理]
\label{thm:general-period-moment-obstruction}
对于每个合法的 \(p\)-周期相位，

\begin{equation}
\begin{aligned}
M_{1}&=4p, \\
M_{2}&=20p+16d, \\
M_{3}&=118p+168d+96a+48b
\end{aligned}
\text{。}
\end{equation}

若 \(R(Q)\leq 8\)，则必有

\begin{equation}
\begin{aligned}
d&\leq \frac{3p}{4}, \\
40d+96a+48b&\leq 42p
\end{aligned}
\text{。}
\end{equation}

这些仅是必要条件。

\end{theorem}
最后，对一个有界但非平凡的周期范围进行分类。

\begin{theorem}[低周期谱极值分类定理]
\label{thm:low-period-frontier}
考虑由式~\eqref{eq:periodic-operator} 定义在无限图上的周期算子，其中 Hamilton 规范
三角形字 \(\tau\) 的最小正周期不超过 16，并以 \(R\) 表示其无限体积谱半径
的平方。
除平移、反射、全局边符号取反以及单位胞重复这些自然等价外，目标相位是唯一
最优相位。重复周期胞仅为证书计数而保留，并通过能区折叠识别为谱等价。

\end{theorem}
该定理是计算机辅助的：一项完整精确的有限分类把闭游走矩压缩与剩余竞争相位的
精确证书结合起来。完整划分和证书计数见第~\ref{sec:low-period-frontier} 节。

\subsection{证明策略}

第一个组织工具是 Hamilton 规范。通过切换，可将所有步长一边变为 \(+1\)，
只留下一个和乐切口；此时一个赋号由三角形通量的周期字 \(\tau\) 编码。其
相邻乘积

\begin{equation}
Q_i=\tau_i \tau_{i+1}\text{，}
\end{equation}

就是四边形通量。周期性允许把算子 Floquet 分解为有限 Hermitian Laurent
矩阵 \(H_Q(z)\)。对目标字，特征行列式关于特征值为偶函数，并约化成变量
\(y=\lambda^{2}\) 与 \(c=z+z^{-1}\) 的四次式 \(P(y,c)\)。在候选谱边处的
一个精确正系数展开证明锐界及其等号情形。

第二个工具是对无限算子作平方。在 \(A_\tau^{2}\) 中，每个奇位移系数都含有
因子 \(1+Q_i\)。因此，负通量恰好消去该耦合，而正通量则引入绝对值为二的
振幅。偶次 Bloch 迹的常数项计数符号闭游走。若所有谱带值的平方均不超过
8，则它们的逐次矩满足 \(M_{k+1}\leq 8M_k\)。因此，正超额量必使谱超过
8。结合这一单向蕴含与缺陷几何，即可证明周期八三分定理并得到一般周期的
必要条件。

\subsection{范围与公开状态}

本文结果区分三个优化区域。定理~\ref{thm:eight-barrier-trichotomy} 讨论表示
周期为八的无限体积相位。定理~\ref{thm:low-period-frontier} 讨论本原 Hamilton
规范周期至多为 16 的周期相位。定理~\ref{thm:smallest-counterexample} 只在已
枚举至 32 阶的范围内讨论全部有限赋号。这些结果均未证明对所有周期、所有
非周期赋号或每个有限阶的全局最优性。

截至 2026 年 8 月 21 日，在源论文、作者列出的后续材料及已记录的专题检索
范围内，尚未发现对猜想 3 的公开直接解决。该文献状态说明仅限于上述检索范围，
不构成绝对优先权主张。

\subsection{文章结构}

第~\ref{sec:preliminaries} 节固定切换、通量、Floquet、矩和等价记号。
第~\ref{sec:smallest-counterexample} 节证明定理~
\ref{thm:smallest-counterexample}。第~\ref{sec:periodic-construction} 节和第~
\ref{sec:period-eight-edge} 节证明周期反例族及其锐谱边。第~
\ref{sec:eight-barrier} 节证明周期八相位三分定理并解释手征对称。第~
\ref{sec:general-period} 节导出一般周期矩和必要条件。第~
\ref{sec:low-period-frontier} 节证明有界低周期谱极值分类。第~
\ref{sec:computational-verification} 节说明计算机辅助证明的边界和可复现性
模型。第~\ref{sec:discussion} 节汇集推论与开放问题。附录给出复核计算机
辅助结论所需的轨道计数论证、精确有限证书和计算协议。
